\documentclass[12pt]{article}
\usepackage[utf8]{inputenc}
\usepackage[T1]{fontenc}
\usepackage{hyperref}
\title{Kemedzung Language Resources for K-12 Curriculum Development}
\author{Compiled by Claude Sonnet 5.5 (claude-sonnet-5-5) for Andr\'as Kornai}
\date{}
\begin{document}
\maketitle

\section{Introduction}
Kemedzung (Kemezung, K\'{\i}m\u{a}dzu\~{n}; ISO 639-3 dmo; Glottocode keme1240) is an Eastern Beboid language of Cameroon's Northwest Region, spoken in the villages of Dumbu and Kwei. Estimates are about 3,450 to 3,540 speakers, with 4,500 given by Global Recordings Network. It is lexically close to Bebe.

\section{Primary Educational Resources}
No primers or school programmes were found in search results. The dictionary below is the main reusable resource.

\section{Linguistic and Reference Materials}
\begin{itemize}
\item Webonary Kemedzung dictionary (about 16,000 words and phrases collected in 2015): \url{https://www.webonary.org/kemedzung/?p=9444}
\item Wikipedia, Kemezung language: \url{https://en.wikipedia.org/wiki/Kemezung_language}
\item Omniglot, Kemezung: \url{https://www.omniglot.com/writing/kemezung.htm}
\item Glottolog reference, Cox 2005 (phonology notes): \url{https://glottolog.org/resource/reference/id/158286}
\end{itemize}

\section{Religious and Community Resources}
\begin{itemize}
\item Scripture Earth page for dmo: \url{https://www.scriptureearth.org/00eng.php?iso=dmo}
\end{itemize}

\section{Digital and Audio Resources}
\begin{itemize}
\item Global Recordings Network, Kemedzung: \url{https://globalrecordings.net/en/language/dmo}
\end{itemize}

\section{Curriculum Development Context}
A published dictionary and orthography notes exist, which gives a base for vocabulary and reader development. No school use was found. Contact details for local organisations were not seen in search results and are not given.
\end{document}